\section{\textit{Machine Learning}}\label{sec:machine_learning}

Dalam bidang \textit{machine learning}, klasifikasi merupakan salah satu tugas utama yang bertujuan untuk memetakan input data ke dalam label kelas tertentu berdasarkan pola yang dipelajari dari data. Secara umum, terdapat empat paradigma utama dalam membangun model pembelajaran mesin, yaitu \textit{unsupervised learning}, \textit{semi-supervised learning}, \textit{supervised learning}, dan \textit{reinforcement learning} (Murphy, 2012; Kelleher et al., 2015; Sutton \& Barto, 2018).

\textit{Unsupervised learning} digunakan ketika data tidak memiliki label kelas, sehingga algoritma berusaha menemukan struktur, pola, atau distribusi tersembunyi dalam data secara mandiri. Sebagai contoh, algoritma \textit{clustering} dapat mengelompokkan pelanggan toko berdasarkan kemiripan kebiasaan belanja mereka tanpa perlu mengetahui kategori pelanggan tersebut sebelumnya. Contoh lain adalah mendeteksi anomali menggunakan algoritma seperti \textit{Isolation Forest}, \textit{Local Outlier Factor}, dan \textit{One-Class SVM} (Chandola, Banerjee, \& Kumar, 2009; Aggarwal, 2013). Teknik ini umum digunakan dalam kasus ketika tidak tersedia informasi eksplisit mengenai kelas data, misalnya untuk memisahkan profil penyewa yang normal dan yang mencurigakan tanpa adanya data riwayat pelanggaran sebelumnya.

\textit{Semi-supervised learning} digunakan pada dataset yang memiliki sedikit data berlabel dan mayoritas data tidak berlabel. Metode ini memanfaatkan sejumlah kecil data berlabel tersebut untuk memandu proses pembelajaran dan mengekstrapolasi label ke data yang belum dilabeli. Pendekatan ini sering digunakan ketika proses pelabelan manual memerlukan waktu dan biaya pakar yang sangat tinggi (Chapelle et al., 2006). Contohnya adalah pengkategorian artikel berita, di mana hanya sebagian kecil artikel yang diberikan label topik (misalnya ``olahraga'' atau ``politik'') oleh manusia, dan algoritma menggunakan pola kata dari artikel tersebut untuk melabeli jutaan artikel lainnya secara otomatis (Zhu, 2005).

Sementara itu, \textit{supervised learning} digunakan ketika seluruh data latih telah dilabeli dengan kelas target tertentu. Model dilatih menggunakan pasangan \textit{input-output} (fitur dan label), sehingga mampu mempelajari pemetaan matematis antar fitur dan menghasilkan prediksi pada data baru (James et al., 2013). Contohnya adalah sistem penyaring email \textit{spam}, di mana model dilatih menggunakan ribuan riwayat email yang sudah ditandai secara eksplisit sebagai ``spam'' atau ``bukan spam'' agar mampu mengklasifikasikan email yang baru masuk. Algoritma yang umum digunakan dalam pendekatan ini meliputi \textit{Decision Tree, Random Forest, Support Vector Machine, K-Nearest Neighbors,} dan \textit{XGBoost} (Tan, Steinbach, \& Kumar, 2019).

Di sisi lain, \textit{reinforcement learning} merupakan pendekatan prediktif, di mana sebuah agen (program) belajar mengambil serangkaian keputusan melalui interaksi langsung dengan lingkungannya untuk memaksimalkan sinyal imbalan (\textit{reward}) (Sutton \& Barto, 2018). Algoritma ini tidak diajari tindakan mana yang benar melalui label data secara statis, melainkan harus mengeksplorasi dan menemukan sendiri tindakan yang menghasilkan \textit{reward} tertinggi melalui mekanisme coba-coba (\textit{trial-and-error}). Contoh penerapannya adalah pada kecerdasan buatan yang belajar bermain catur; agen tersebut akan mendapat imbalan (+1) saat memakan bidak lawan, dan mendapat hukuman atau penalti (-1) saat membuat langkah yang merugikan (Silver et al., 2016).

Dalam proses pelatihan model \textit{machine learning}, salah satu tantangan terbesar yang sering terjadi adalah \textit{overfitting}. \textit{Overfitting} merupakan kondisi di mana model menghafal data pelatihan dengan terlalu detail, termasuk menangkap gangguan kecil (\textit{noise}) dan variasi acak yang tidak penting (Dietterich, 1995). Akibatnya, model tampak memiliki tingkat akurasi yang sangat luar biasa tinggi pada data latih, namun ketika diuji menggunakan data baru yang belum pernah dilihat sebelumnya, kemampuannya menurun secara drastis (Hawkins, 2004). Analogi sederhananya seperti seorang siswa yang hanya menghafal soal-soal dan jawaban dari buku latihan ujian tanpa benar-benar memahami konsep dasar materinya. Saat menghadapi ujian sebenarnya dengan angka yang sedikit diubah, ia akan gagal menjawab (Goodfellow, Bengio, \& Courville, 2016). Untuk mengatasi masalah \textit{overfitting}, beberapa strategi umum yang diterapkan meliputi penggunaan lebih banyak data pelatihan, membatasi kerumitan model, ataupun menggunakan metode regularisasi (Hastie, Tibshirani, \& Friedman, 2009).

% Berdasarkan karakteristik data yang telah dilabeli sebagai normal atau tidak normal, maka pendekatan yang digunakan dalam penelitian ini adalah \textit{supervised learning}. Hal ini memungkinkan pelatihan model klasifikasi yang dapat mengenali pola perilaku berdasarkan data historis, serta mengevaluasi performanya secara objektif menggunakan metrik evaluasi seperti akurasi, presisi, dan recall.



